\rok{Фебруар 2024.}{А - Додатни поправни тест другог теста}

Додатни поправни тест из Алгоритама и структура података

Други практични тест одржан 06.02.2024.

\zadatak{1}{Имплементација DFSVisit методе Depth-First-Search алгоритма}
Написати имплементацију \texttt{DFSVisit} методе истоименог алгоритма.

\textbf{Додатне напомене за решавање задатка:}
\begin{itemize}
	\item У template-у је закоментарисана метода:
	\begin{Verbatim}[fontsize=\footnotesize, numbers=left, xleftmargin=10mm]
private void DFSVisit(LinkedList.Node node, int nodeIndex)
	\end{Verbatim}
	\item Поменута метода је основна метода за рекурзивно извршавање DFS алгоритма.
	\item У Main методи обезбедити начин за тестирање и проверу нове функционалности.
\end{itemize}



\zadatak{2}{Замена понављајућих и непонављајућих карактера}
Написати алгоритам временске комплексности $O(n)$ који ће из string-а који се прослеђује раздвајати понављајуће од непонављајућих карактера и замењивати их према следећем шаблону:

\begin{itemize}
	\item Понављајући карактер = \texttt{@}
	\item Непонављајући карактер = \texttt{\#}
\end{itemize}

\textbf{Додатне напомене за решавање задатка:}
\begin{itemize}
	\item Користити Hash табелу која је дата у шаблону.
	\item Имплементацију писати у оквиру закоментарисане методе:
	\begin{Verbatim}[fontsize=\footnotesize, numbers=left, xleftmargin=10mm]
public static string changeAllDuplicates(string str)
	\end{Verbatim}
	\item Алгоритам реализовати тако да не постоји разлика између малих и великих карактера.
	\item За конверзију string-а у низ карактера користити \texttt{ToCharArray()} по принципу: \texttt{str.ToCharArray()}.
	\item Карактере није потребно експлицитно конвертовати у целобројни тип приликом убацивања у Hash табелу, али приликом каснијег исписивања користити експлицитно кастовање \texttt{(char)}.
	\item Редослед карактера приликом исписивања није битан.
\end{itemize}

\textbf{Примери:}
\begin{itemize}
	\item \textbf{Улаз 1:} \texttt{"Algoritam"}
	\item \textbf{Излаз 1:} \texttt{"@######@#"} (само карактер \texttt{A} се појављује више пута)
	\item \textbf{Улаз 2:} \texttt{"Strukture podataka"}
	\item \textbf{Излаз 2:} \texttt{"#@@@@@@@# ###@@@@"} (карактери \texttt{a}, \texttt{r}, \texttt{u}, \texttt{k} и \texttt{t} се појављују више пута)
\end{itemize}



\zadatak{3}{Проналажење разлике чворова на парним и непарним нивоима стабла}
У оквиру стандардне структуре класе \texttt{BinarySearchTree} написати имплементацију методе која враћа вредност добијену одузимањем вредности чворова на непарним нивоима од вредности чворова на парним нивоима стабла.

\textbf{Додатне напомене за решавање задатка:}
\begin{itemize}
	\item Корен стабла се налази на нултом нивоу.
	\item У оквиру класе \texttt{BinarySearchTree} креирати јавну методу са потписом:
	\begin{Verbatim}[fontsize=\footnotesize, numbers=left, xleftmargin=10mm]
public int findDiffEvensOdds()
	\end{Verbatim}
	\item Из претходно креиране јавне методе треба да се позива приватна метода са потписом:
	\begin{Verbatim}[fontsize=\footnotesize, numbers=left, xleftmargin=10mm]
private int findDiffEvensOdds(Node curr, bool levelIsEven)
	\end{Verbatim}
	\item Приватна метода треба да садржи имплементацију алгоритма за рачунање разлике вредности чворова на парним и непарним нивоима стабла.
	\item Алгоритам \textbf{ОБАВЕЗНО} писати у оквиру класе \texttt{BinarySearchTree}.
\end{itemize}

\textbf{Примери:}

\begin{center}
\begin{minipage}{\textwidth}
\centering
\begin{forest}
for tree={
	circle,
	draw,
	thick,
	fill=gray!20,
	minimum size=8mm,
	inner sep=1pt,
	s sep=8mm,
	l sep=12mm,
	edge={-, thick},
	font=\small
}
[8
	[4
		[2
			[, phantom]
			[3]
		]
		[5]
	]
	[9
		[, phantom]
		[11]
	]
]
\end{forest}

\vspace{0.3cm}
\textbf{Слика 1.} Пример BST-а.
\end{minipage}
\end{center}

\begin{itemize}
	\item \textbf{Улаз 1:} BST са слике 1
	\item \textbf{Излаз 1:} 10
\end{itemize}

\begin{center}
\begin{minipage}{\textwidth}
\centering
\begin{forest}
for tree={
	circle,
	draw,
	thick,
	fill=gray!20,
	minimum size=8mm,
	inner sep=1pt,
	s sep=8mm,
	l sep=12mm,
	edge={-, thick},
	font=\small
}
[10
	[4
		[2
			[3]
			[, phantom]
		]
		[6
			[5]
			[9]
		]
	]
	[20
		[12
			[, phantom]
			[14]
		]
		[29
			[27]
			[, phantom]
		]
	]
]
\end{forest}

\vspace{0.3cm}
\textbf{Слика 2.} Пример BST-а.
\end{minipage}
\end{center}

\begin{itemize}
	\item \textbf{Улаз 2:} BST са слике 2
	\item \textbf{Излаз 2:} -38
\end{itemize}



\sablon{Шаблон додатног поправног теста фебруар 2024. група А}{2023/Februar/GrupaA/AISP2023_PopravniFebruarDodatniTermin_Test2_GrupaA.linq}
